% Numerical setup.  Written 2026-09-07 from the code as it stands
% (HLLD @ 43bc6c2, rmhd_final @ a2451c7).  Every number here is taken from
% the source, not from memory; see docs/numerical_setup_notes.md for the
% file and line each statement comes from.
\section{Numerical Setup}
\label{sec:setup}

Two codes are used, and they answer different questions. The first is a
special-relativistic magnetohydrodynamics (SRMHD) code that evolves the fluid
on a grid; the second is an exact Riemann solver, which the first can call as
its flux function. The exact solver is also used on its own, as the reference
against which the approximate fluxes are measured.

\subsection{The evolution code}
\label{sec:setup:evolution}

The fluid is evolved with a finite-volume Godunov scheme in conservative
form. The evolved variables are the conserved density $D$, the momentum
density $S_i$, the energy density $\tau$, and, in two dimensions, the
out-of-plane field $B^z$; the in-plane field is carried on the cell faces
(Sect.~\ref{sec:setup:ct}). We close the system with an ideal-gas equation of
state, $p = (\Gamma - 1)\rho\epsilon$, with $\Gamma = 5/3$ for the rotor and
$\Gamma = 4/3$ for the Orszag--Tang vortex.

\paragraph{Reconstruction.}
The primitive variables $(\rho, p, v^i, B^i)$ are reconstructed to the cell
interfaces by a piecewise-linear method (PLM). The cell slope is the
monotonised-central limit of the two one-sided differences,
%
\begin{equation}
  \sigma_i = \mathrm{mc}\bigl(Q_i - Q_{i-1},\; Q_{i+1} - Q_i\bigr),
  \qquad
  \mathrm{mc}(a,b) =
  \begin{cases}
    \mathrm{sgn}(c)\,\min\bigl(2\lvert a\rvert,\, 2\lvert b\rvert,\,
                               \lvert c\rvert\bigr), & ab > 0,\\[2pt]
    0, & \text{otherwise},
  \end{cases}
  \label{eq:mc}
\end{equation}
%
with $c = (a+b)/2$, and the interface states are $Q_i \pm \sigma_i/2$. The
method is second order where the solution is smooth and reduces to first order
at extrema, and it needs two ghost zones. The minmod limiter and
piecewise-constant reconstruction are available and are used for comparison.

Two departures from a naive variable-by-variable PLM are needed in the
relativistic case. First, the reconstructed velocity variable is the spatial
part of the four-velocity, $z^i = W v^i$, from which
$v^i = z^i / \sqrt{1 + \lvert \mathbf{z} \rvert^2}$ is recovered; this makes
$\lvert \mathbf{v} \rvert < 1$ true at every interface by construction, whereas
limiting the three components of $v^i$ independently can produce a superluminal
interface state wherever the limiter is active near a steep velocity gradient.
That is not a hypothetical failure for the rotor, whose initial data already
reach $W \simeq 10$. Second, the specific internal energy is recomputed from
the reconstructed $(p, \rho)$ rather than reconstructed on its own, so that the
interface state lies on the equation-of-state surface. Finally, on any face
where the limited state would violate the density or pressure floor, that face
falls back to piecewise-constant reconstruction.

\paragraph{Riemann solver.}
The interface flux is computed by one of four solvers, selected at run time:
HLLE, HLLC, HLLD, or the exact solver of Sect.~\ref{sec:setup:exact}. This is
the only difference between the runs compared in this work; the grid, the
reconstruction, the time integrator, the constrained transport and the
conservative-to-primitive inversion are held fixed, so that any difference in
the solution is attributable to the flux function alone.

\paragraph{Time integration.}
Time is advanced with the strong-stability-preserving Runge--Kutta method of
Shu \& Osher, third order by default and second order optionally. The time
step is set by the unsplit Courant condition on the fast magnetosonic speed
in both directions, with a Courant factor of $0.25$.

\paragraph{Conservative-to-primitive inversion.}
The primitive variables are recovered with the vectorised scheme of
Kastaun et al., which reduces the inversion to a one-dimensional root find of
a master function with a bracketed root, and is therefore robust in the
magnetically dominated and highly relativistic regimes that the rotor reaches.
The inversion is applied to every cell after every Runge--Kutta stage.

\paragraph{Constrained transport.}
\label{sec:setup:ct}
In two dimensions the magnetic field is kept solenoidal by constrained
transport: the in-plane field is stored on cell faces, and the corner
electromotive force is built from the interface fluxes by the upwinded
averaging of Gardiner \& Stone. The cell-centred and face-centred variables
are combined with identical Runge--Kutta coefficients at every stage, which is
what keeps each stage divergence-free without dedicated machinery. The
resulting $\lvert\nabla\cdot\mathbf{B}\rvert$ remains at the level of
round-off, and is reported with the results.

\subsection{The exact Riemann solver}
\label{sec:setup:exact}

The second code solves the SRMHD Riemann problem exactly, in the sense that
every wave is represented by its own jump or integral relation rather than by
an approximate average. We follow the formulation of Giacomazzo \& Rezzolla
for the case $B^x \neq 0$, in which the solution consists of seven waves
separating eight uniform states: two fast waves, two rotational (Alfv\'en)
discontinuities, two slow waves, and the contact discontinuity.

\paragraph{Extension to six unknowns.}
The original formulation carries four unknowns, the total pressure behind each
fast wave and the two tangential field components at the contact, and closes
the system by requiring the normal and tangential velocities and the total
pressure to be continuous across the contact. The rotational discontinuities
are then left to an inner iteration with case-dependent starting values. That
inner iteration is the weak point: a rotational discontinuity admits more than
one branch, and the outer Newton has no control over which branch the inner
solve selects, so the outer iteration can converge to a structure that is not
the solution of the problem posed.

We therefore promote the two rotation angles to unknowns of the outer system.
The unknown vector becomes
%
\begin{equation}
  \mathbf{u} = \bigl[\,\ln p^\ast_{\mathrm{LF}},\;
                      \ln \lvert \mathbf{B}_t^\ast \rvert,\;
                      \psi^\ast,\;
                      \ln p^\ast_{\mathrm{RF}},\;
                      \varphi_L,\; \varphi_R \,\bigr],
  \label{eq:unk6}
\end{equation}
%
where $p^\ast_{\mathrm{LF}}$ and $p^\ast_{\mathrm{RF}}$ are the total
pressures behind the left and right fast waves, $\lvert\mathbf{B}_t^\ast\rvert$
and $\psi^\ast$ are the magnitude and orientation of the tangential field at
the contact, and $\varphi_L$, $\varphi_R$ are the angles through which each
rotational discontinuity turns the tangential field. Writing the magnitudes as
logarithms and the orientations as angles makes every unknown of order unity
and positivity automatic, so a single absolute step size is well conditioned
for the whole vector.

Each Alfv\'en wave is then solved with its rotation \emph{pinned} to the
corresponding angle unknown, which makes the wave deterministic and moves the
branch choice into the outer Newton, where it belongs. The system is closed by
six residuals: continuity of the three velocity components and of the total
pressure across the contact,
%
\begin{equation}
  \llbracket v^x \rrbracket = \llbracket v^y \rrbracket =
  \llbracket v^z \rrbracket = \llbracket p_{\rm tot} \rrbracket = 0 ,
  \label{eq:cd}
\end{equation}
%
together with one consistency condition per slow wave. The slow waves are
parameterised by the tangential field they must produce at the contact, which
in general over-determines them; the slack is the residual of the transverse
momentum jump condition for a slow shock, and the mismatch in field
orientation for a slow rarefaction.

\paragraph{Wave solvers.}
Each wave is solved by its own routine. Fast and slow shocks are obtained from
the Rankine--Hugoniot conditions by a pressure-parameterised root find over
the admissible speed window, bracketed by a scan of the residual so that the
correct branch is selected rather than the nearest root. Rarefactions are
obtained by integrating the self-similar ordinary differential equations
through the fan with a fourth-order Runge--Kutta scheme. Rotational
discontinuities are solved for the pinned rotation angle. Where a wave family
degenerates, for instance where the Alfv\'en speed coincides with a
magnetosonic speed, the corresponding reduced three-wave problem is solved
instead; the structure of each interface is classified from its eigenvalues
before the solve.

\paragraph{Outer iteration.}
The six-dimensional system is solved by a damped Newton method with a
forward-difference Jacobian. The step is taken through a singular-value
decomposition with truncation of directions below the finite-difference noise
level, which is necessary because the Jacobian is genuinely rank-deficient
whenever the problem is coplanar and the rotation angles are unobservable. A
trust region caps the step, and a backtracking line search accepts a step only
if it reduces the residual; the best iterate is retained. Convergence is
declared at $\lVert \mathbf{f} \rVert_\infty < 10^{-8}$.

\paragraph{Initial guess and multistart.}
The Newton iteration converges only from a starting point inside the basin of
the solution, and that basin shrinks as the magnetisation grows. The starting
point is therefore supplied by a neural network, a feed-forward
multilayer perceptron trained on forward-constructed solutions: it maps
thirteen scale- and rotation-invariant features of the state pair to the nine
numbers from which Eq.~(\ref{eq:unk6}) is assembled, with the three angles
represented by their sine and cosine so that the training loss is continuous
across the branch cut. Interfaces that fail to converge from the first guess
are retried from further starting points, drawn either from other trained
checkpoints or by perturbing the first. Retries recover a substantial fraction
of the failures, which is the expected behaviour when the failures are basin
misses rather than an insufficient iteration budget.

\paragraph{Use as a flux function and fallback.}
When the exact solver provides the flux for a grid run, all interfaces of a
sweep are solved in one batched call. An interface is given the exact flux
only if the Newton converged, the state on the ray $\xi = x/t = 0$ could be
resolved, that state is physical, and the structure behind it satisfies the
complete seven-wave jump conditions to $10^{-8}$; otherwise the interface
keeps the HLLD flux. The number of interfaces attempted and the number
verified exact are reported with every run.

Two classes of interface are not attempted. The first is those whose left and
right states differ by less than a relative threshold $\tau_{\rm weak} =
10^{-2}$, measured by a scale-free rotation-invariant jump that compares
$\rho$ and $p$ relatively and $\mathbf{v}$ and $\mathbf{B}$ as vector
differences. This is where the choice of solver cannot matter: taking real
interfaces and shrinking the discontinuity about the mean state, the relative
difference between the exact and HLLD fluxes falls very nearly linearly with
the jump, reaching a median of $2.5 \times 10^{-5}$ at the threshold and
$8 \times 10^{-8}$ two decades below it, with no floor. The gate excludes
$85.9\%$ of interfaces at a cost estimated at $0.02\%$ of the solution.

The second class is interfaces whose fan lies entirely on one side of the
ray. HLLD's signal speeds are clamped at zero, so a vanishing left speed
makes its flux exactly $\mathbf{F}_L$, which is also the exact solution
there. Measured over $45\,650$ interfaces this branch fires on $29\%$ of
those attempted, and on every one of them the two fluxes agree to at worst
$1.9 \times 10^{-16}$, while on the complement not one agrees to better than
$10^{-14}$. Skipping them costs $1.59$ times less work in the solver stage
and changes no answer.
